\documentclass[withoutpreface]{cumcmthesis}

% 严格继承用户模板；只增加当前题目确有需要的宏包。
\usepackage[framemethod=TikZ]{mdframed}
\usepackage{url}
\usepackage{subcaption}
\usepackage{tikz}
\usetikzlibrary{arrows.meta,positioning,calc,fit,shapes.geometric}
\usepackage{mathtools}

\title{论文题目}
\tihao{C}
\baominghao{xxxx}
\schoolname{XX大学}
\membera{ }
\memberb{ }
\memberc{ }
\supervisor{ }
\yearinput{2026}
\monthinput{09}
\dayinput{ }

\begin{document}
\maketitle

\begin{abstract}
% 第一段：背景矛盾 + 数据/机制特点 + 全文统一目标，2--4句。

针对问题一，考虑到……，因此建立\textbf{……模型}，通过……得到……，并利用……验证……。

针对问题二，在问题一……的基础上，为解决……，建立\textbf{……模型}。结果表明……。

针对问题三，……。

针对问题四，……。

\keywords{关键词1\quad 关键词2\quad 关键词3\quad 关键词4}
\end{abstract}
\newpage

\section{问题重述}
\subsection{问题背景}
\subsection{需要解决的问题}

\section{问题分析}
\subsection{问题一的分析}
\subsection{问题二的分析}
\subsection{问题三的分析}
\subsection{问题四的分析}

% 在问题分析末尾放高质量总览图，展示数据、模型、结果、验证的传递关系。

\section{模型假设}

\section{符号说明}
\begin{table}[!htbp]
\caption{主要符号说明}\label{tab:symbols}\centering
\begin{tabular}{ccc}
\toprule[1.5pt]
符号 & 含义 & 单位\\
\midrule[1pt]
$x$ & 示例变量 & --\\
\bottomrule[1.5pt]
\end{tabular}
\end{table}

\section{模型的建立与求解}

\subsection{问题一的模型建立与求解}
\subsubsection{数据特征与建模动机}
\subsubsection{模型建立}
\subsubsection{模型求解与结果分析}
\subsubsection{模型验证}

\subsection{问题二的模型建立与求解}

% 重要优化模型在推导后集中总括，例如：
% \begin{equation}
% \boxed{
% \left\{
% \begin{aligned}
% \max_{\boldsymbol{x}}\quad & F(\boldsymbol{x}),\\
% \text{s.t.}\quad & g_i(\boldsymbol{x})\le 0,\\
% & h_j(\boldsymbol{x})=0,\\
% & \boldsymbol{x}\in\Omega.
% \end{aligned}
% \right.
% }
% \end{equation}

\subsection{问题三的模型建立与求解}

\subsection{问题四的模型建立与求解}

\section{模型评价与推广}

\begin{thebibliography}{99}
\bibitem{ref1} 作者. 文献标题[J]. 期刊, 年, 卷(期): 页码.
\end{thebibliography}

\begin{appendices}
\section{支撑材料文件清单}
\begin{itemize}
\item code/：完整可运行代码；
\item figures/：正文图；
\item data/：自主查阅/处理的数据；
\item support/：较大中间结果及按当年规则需要的其他支撑材料。
\end{itemize}
\end{appendices}

\end{document}
